Weight averaging offers a way to evaluate a training trajectory through more
than its final iterate. Diffusion-model implementations often retain an EMA of
the trained parameters, including the original DDPM work \cite{Ho2020}. This
practice does not establish that every averaging rule improves every endpoint.
The optimizer, learning-rate path, averaging window and evaluation procedure
all affect the comparison.

A smaller difference between averaged and raw weights also has two possible
interpretations: the average may have become worse, or the raw model may have
caught up. A useful comparison therefore reports both the paired relative
effect and the absolute quality of each model. The distinction matters when
comparing constant learning rates with a schedule that approaches zero.

The retained study asks how the relative benefit of fixed-decay EMA changes
across terminal training duration and the complete learning-rate policy in a
small DDPM. It separates complete cosine schedules at two durations from a
constant-rate trajectory evaluated at the same update counts. The study uses
simple two-dimensional distributions so that paired model states, generated
samples and inputs can be retained and checked within a finite local budget.
It does not seek a new state of the art or establish a general explanation of
averaging. In particular, the relationship between averaging and annealing is
already directly examined by Ajroldi et al.\ \cite{Ajroldi2025}.

This paper is a publication revision of completed Allagma run r07. The
literature search described here occurred after that experiment. Its role is to
correctly situate and qualify the retained findings, not to retroactively change
the protocol or claim that these sources determined its original freeze.
